\documentclass[12pt,letterpaper]{article}
\usepackage{amsmath}
\usepackage{amsthm}
\usepackage{amsfonts}
\usepackage{amssymb}
\usepackage{amscd}
\usepackage{enumerate}
\usepackage{fancyhdr}
\usepackage{mathrsfs}
\usepackage{bbm}
\usepackage{framed}
\usepackage{mdframed}
\usepackage{cancel}
\usepackage{mathtools}
\usepackage{verbatim}
\usepackage[letterpaper,voffset=-.5in,bmargin=3cm,footskip=1cm]{geometry}
\setlength{\parindent}{0.0in}
\setlength{\parskip}{0.1in}
\allowdisplaybreaks
\headheight 15pt
\headsep 10pt
\newcommand\N{\mathbb N}
\newcommand\Z{\mathbb Z}
\newcommand\R{\mathbb R}
\newcommand\Q{\mathbb Q}
\newcommand\lcm{\operatorname{lcm}}
\newcommand\setbuilder[2]{\ensuremath{\left\{#1\;\middle|\;#2\right\}}}
\newcommand\E{\operatorname{E}}
\newcommand\V{\operatorname{V}}
\newcommand\Pow{\ensuremath{\operatorname{\mathcal{P}}}}

\DeclarePairedDelimiter\ceil{\lceil}{\rceil}
\DeclarePairedDelimiter\floor{\lfloor}{\rfloor}
\newcommand\hint[1]{\textbf{Hint}: #1}
\newcommand\note[1]{\textbf{Note}: #1}
\newcounter{enum22i}
\newenvironment{22enumerate}{\begin{enumerate}[a.]\itemsep0em\setcounter{enumi}{\value{enum22i}}}{\setcounter{enum22i}{\value{enumi}}\end{enumerate}}
\newenvironment{22itemize}{\begin{itemize}\itemsep0em}{\end{itemize}}
\fancypagestyle{firstpagestyle} {
  \renewcommand{\headrulewidth}{0pt}
  \lhead{\textbf{CSCI 0220}}
  \chead{\textbf{Discrete Structures and Probability}}
  \rhead{M. Littman}
}

\fancypagestyle{fancyplain} {
  \renewcommand{\headrulewidth}{0pt}
  \lhead{\textbf{CSCI 0220}}
  \chead{Midterm Practice}
}
\pagestyle{fancyplain}
\usepackage{tikz}
\usepackage{graphicx}

\begin{document}
  \thispagestyle{firstpagestyle}
  \begin{center}
    {\huge \textbf{Midterm Practice}}
  \end{center}
  
\section{Proofs}

\subsection*{Problem 1}
{
\nopagebreak 

{\setcounter{enum22i}{0}
  Show that the square of an even number is also an even number using both:
  \begin{22itemize}
    \item A Direct Proof
    \item A Proof by Contradiction
  \end{22itemize}

}
}

\subsection*{Problem 2}
{
\nopagebreak 

{\setcounter{enum22i}{0}
  Prove that $\forall n \in \Z,\ n^2 + 3n - 5$ is odd by dividing into cases.

}
}

\newpage

\section{Logic/Circuits}


\subsection*{Problem 1}
{
\nopagebreak 

{\setcounter{enum22i}{0}

      \begin{22enumerate}
        \item Suppose $p$ and $q$ are propositions such that
        $p$ IMPLIES $q$ is False.  Determine the values of:
          \begin{enumerate}[(i)]\itemsep0pt
            \item (NOT $p$) IMPLIES $q$
            \item $p$ OR $q$
            \item $q$ IMPLIES $p$
          \end{enumerate}
         \item You are now going to design a circuit that takes as
	 input two 1-bit binary numbers
        $A$ and $B$ and outputs whether or not $A > B$, $A < B$,
       	or $A=B$. Namely, the circuit should have two inputs, the bits
        $A$ and $B$, and three outputs $G$, $E$, and $L$.
        For any input, exactly one of $G$ (greater), $E$ (equal),
	or $L$ (less) should be 1. Specifically, $G$ corresponds to
	$A>B$.
          \begin{enumerate}[(i)]\itemsep0pt
            \item Write out a truth table equivalent to this circuit.
            \item Draw your circuit. Please use only \textsc{And},
            \textsc{Or}, and \textsc{Not} gates with at most
            two inputs per gate. Note that we care only about the
	    correctness of the circuit, not its complexity. Be sure to explain how your circuit works.
          \end{enumerate}
      \end{22enumerate}
    
}
}

\subsection*{Problem 2}
{
\nopagebreak 

{\setcounter{enum22i}{0}

      \begin{enumerate}[a.]
        \item Add parentheses to the following expressions to make them true (1 represents
           true, and 0 represents false). Note also that there is no implicit
           ordering; that is, all ordering comes from your parentheses. State
           explicitly any assumptions you are making about the order of operations.
          \begin{enumerate}[i.]
\renewcommand{\land}{\mathrm {\ AND\ }}
\renewcommand{\lor}{\mathrm {\ OR\ }}
\renewcommand{\Rightarrow}{\mathrm {\ IMPLIES\ }}
\renewcommand{\Leftrightarrow}{\mathrm {\ IFF\ }}

          \item $0 \land 1 \lor 1 \Rightarrow 1 \land 1 \land 1 \lor 0$
          \item $0 \lor 0 \land 1 \land 0 \land 1 \land 1 \lor 1$
          \item $0 \land 1 \lor 1 \Leftrightarrow 0 \Rightarrow 0$
          \item $0 \lor 1 \Rightarrow 0 \land 0 \Rightarrow 1$
          \end{enumerate}
	\item Prove the following logical equivalences:
	
\renewcommand{\land}{\mathrm {\ AND\ }}
\renewcommand{\lor}{\mathrm {\ OR\ }}
\renewcommand{\Rightarrow}{\mathrm {\ IMPLIES\ }}
\renewcommand{\Leftrightarrow}{\mathrm {\ IFF\ }}
\renewcommand{\lnot}{\mathrm {NOT\ }}
	  \begin{enumerate}[i.]	  
	  
		\item $\lnot (p \lor (\lnot p \land q) \equiv \lnot p \land \lnot q$
		\item $(p \Rightarrow r) \land (q \Rightarrow r) \equiv (p \lor q) \Rightarrow r$
	  \end{enumerate}
	\item Prove that $q \land \lnot (p \Rightarrow q)$ is unsatisfiable.
	\item Prove that $(p \land q) \Rightarrow (p \lor q)$ is valid.
	\end{enumerate}
    
}
}



\newpage

\section{Set Theory}

\subsection*{Problem 1}
{
\nopagebreak 

{\setcounter{enum22i}{0}
	
    Determine whether each of the following statements is true or false, and explain why.
      \begin{22enumerate}
      \item The powerset of the empty set has the empty set as a member.
      \item The empty set is an element of every set.
      \item The empty set is a subset of any set that does not have the empty set as a member.
      \item Any set that has the empty set as a member must be the empty set.
      \item The set containing the set containing the set containing the empty set has a cardinality of zero (here, $A$ contains $B$ means $B \in A$).
      \item The set of all empty sets is the same as the powerset of the empty set.
    \end{22enumerate}
    
    
}
}

\subsection*{Problem 2}
{
\nopagebreak 

{\setcounter{enum22i}{0}


     \begin{22enumerate}
     \item Disprove the following claim: 

     For any two sets finite $A$ and $B$, there are the same number of elements in $\Pow(A \times B)$ and $\Pow(A) \times \Pow(B)$.


     \item Consider two finite, \textit{non-empty} sets $A$ and $B$. Suppose $\Pow(A \times B)$ and $\Pow(A) \times \Pow(B)$ were equal in size. What must be the size of $A$? What must be size of $B$? Explain your answer. 

     \end{22enumerate}
    
}
}

\newpage

\section{Relations}

\subsection*{Problem 1}
{
\nopagebreak 

{\setcounter{enum22i}{0}

    \begin{22enumerate}
      \item
        Prove or disprove that a relation that is reflexive and symmetric is necessarily
        transitive.
      \item
        Prove or disprove that a relation that is reflexive and transitive is necessarily
        symmetric.
      \item
        Prove or disprove that a relation that is symmetric and transitive is necessarily
         reflexive.
    \end{22enumerate}
    
}
}



\subsection*{Problem 2}
{

{\setcounter{enum22i}{0}
	Let $S=\{0,1\}, \; T = \{t | t \subseteq S \times S\}$, and $R$ be the set of all possible functions from $S$ to $S$.
	\begin{22enumerate}
	\item Can an injection from $T$ to $R$ exist?
	If so, give one such injection and prove that this mapping is indeed
	injective.  If not, prove why such a mapping cannot exist.
	\item Can a surjection from $T$ to $R$ exist?
	If so, give one such surjection and prove that this mapping is indeed
       	surjective.  If not, prove why such a mapping cannot exist.
	\item Can a bijection from $T$ to $R$ exist? If so, why?
	If not, why not?
	\end{22enumerate}

}
}


\newpage

\section{Induction}

\subsection*{Problem 1}
{
\nopagebreak 

{\setcounter{enum22i}{0}

        Suppose we have a sequence defined as follows:
    	\begin{align*}
    		a_1 &= 1 \\
    		a_2 &= 3 \\
    		a_k &= a_{k-2} + 2a_{k-1}
    	\end{align*}

    	Prove by induction that $a_n$ is odd for all positive integers $n$.
    
}
}


\subsection*{Problem 2}
{

{\setcounter{enum22i}{0}
	Use mathematical induction to prove the following generalization of one of DeMorgan's laws:
			\[\overline{\bigcap_{j=1}^{n}{A_j}} = \bigcup_{j=1}^{n}{\overline{A_j}},\]
	whenever $A_1,A_2,...,A_n$ are subsets of a universal set $U$ and $n \geq 2$. To give you an idea of what this notation means, note that:
\[\bigcap_{j=1}^{3}{A_j} = A_1 \cap A_2 \cap A_3. \]
  \textbf{Note:} You must prove DeMorgan's Law for the base case.

}
}

\end{document}